\chapter{Trabajo a futuro}
\label{chapter:7-Trabajo a futuro}
En esta sección se detallarán algunos de los avances necesarios para mejorar el proyecto actual. 

Para cumplir nuestra visión y de que los estudiantes utilicen esta herramienta en el desarrollo de sus cursos, es necesario integrar el proyecto dentro de un entorno institucional, mantenido por o con el apoyo directo de la URI (Unidad de Repositorios Institucionales) y OpenFing. La integración garantizaría la sostenibilidad, estabilidad y el respaldo necesario a largo plazo. Durante el transcurso del proyecto, el sistema estuvo alojado en la plataforma "Mi Nube" de Antel como solución provisional mientras se avanza hacia la implementación de la infraestructura definitiva. Una posibilidad es evaluar si la URI pudiera asumir la responsabilidad del despliegue y mantenimiento del entorno operativo para este sistema. Otro punto clave también es sería rediseño del actual procedimiento para subir material a OpenFing para que las clases nuevas se incorporen automáticamente al sistema. De esta manera se puede mantener el material actualizado a lo largo del tiempo sin necesidad de un trabajo manual. 

Ya cubierto el contenido en OpenFing, aun es posible aumentar el contenido para que otras facultades logren integrarse en el sistema. Las otras facultades de la  UDELAR están comenzando a subir sus clases de manera online. Desarrollando herramientas similares a OpenFing. Durante nuestra exposición en Ingeniería de Muestra (ver \autoref{chapter:A-0-Anexo I-Ingeniería de Muestra}), se nos comentó lo útil que puede ser este proyecto para otras facultades. Cubrir otras facultades requeriría un esfuerzo acotado. Se necesitaría extender el pipeline de procesamiento para incorporar clases de otras fuentes y añadir un filtrado por facultad de origen.

Este filtrado resultaría esencial en caso de añadir clases de distintas facultades debido a que muchos temas son abordados en cursos de diferentes facultades. Sin embargo, sí se realizan distinciones con respecto al enfoque, nivel de profundidad y contextos académicos. Por lo que permitir segmentar la información por facultad, se garantiza que el chatbot acceda únicamente a las fuentes más pertinentes para cada caso. De esta forma, las respuestas entregadas no solo serán más precisas y relevantes, sino que también mantendrán coherencia con la perspectiva institucional y disciplinar correspondiente. Esto contribuye a una experiencia más clara, confiable y adaptada a las necesidades específicas de cada estudiante o usuario.
 
Otro trabajo a futuro sería evaluar distintas metodologías y técnicas sobre el sistemas. Por ejemplo, una de estas técnicas a evaluar puede ser la estrategias de \textit{Re-Ranking} sobre los vectores recuperados. Esta estrategia es una modificación dentro del contexto RAG. La arquitectura RAG depende en gran medida de la calidad y relevancia del contexto recuperado de la base de datos vectorial. Dentro del contexto RAG tradicional, los documentos son recuperados en función de similitud matemática entre vectores, independiente de la similitud semántica que los vectores puedan tener. Dependiendo el caso, puede ocurrir que los vectores matemáticamente más similares, no sean semánticamente relevantes, resultando en vectores de poco valor para responder la consulta del usuario.  El \textit{Re-Ranking} es una estrategia que busca refinar este proceso mediante un reordenamiento posterior de los documentos recuperados, priorizando aquellos que presentan una mayor relevancia semántica y contextual, segmentos que prodrían resultar más utiles para la generación de respuestas. En \cite{irany2026hybrid}, los autores implementan una estrategia híbrida de recuperación y \textit{Re-Ranking} en un sistema RAG orientado al dominio biomédico, obteniendo mejoras en la calidad del contexto recuperado y alcanzando resultados favorables en tareas de generación fundamentada en evidencia.

En una línea complementaria, la recuperación actual se basa exclusivamente en similitud vectorial densa, lo que captura bien la similitud semántica entre la consulta y los fragmentos indexados, pero puede no rendir de forma óptima ante términos técnicos exactos, siglas o nombres de asignaturas poco frecuentes en el espacio vectorial. Una estrategia de búsqueda híbrida que combine la recuperación densa actual con un componente de recuperación léxica dispersa, como BM25, podría mejorar la precisión en este tipo de consultas. En ese esquema, los fragmentos recuperados por ambos métodos se fusionarían antes de la etapa de generación, aprovechando las fortalezas complementarias de cada enfoque.

Una mejora transversal que habilitaría varias de las anteriores es la incorporación de un mecanismo de retroalimentación explícita del usuario. Agregar controles simples como valoraciones de calidad sobre las respuestas del asistente permitiría acumular datos reales sobre qué respuestas son consideradas útiles y cuáles no. Este conjunto de interacciones etiquetadas no solo facilitaría el monitoreo continuo de la calidad del sistema en producción, sino que constituiría la materia prima necesaria para procesos de ajuste fino posteriores.

Otro trabajo a futuro relevante es el ajuste fino (\textit{fine-tuning}) del modelo de lenguaje sobre el dominio académico uruguayo. En una arquitectura RAG, la etapa de recuperación determina qué fragmentos del corpus llegan al modelo, pero es la etapa de generación la que decide cómo sintetizarlos en una respuesta. El modelo actual, Llama 3.1 8B, opera sin adaptación de dominio específica, por lo que puede producir respuestas que mezclan el contexto recuperado con conocimiento general del preentrenamiento, lo que aumenta el riesgo de inexactitudes en temas muy específicos de la FING. Ajustar el modelo sobre interacciones reales etiquetadas como correctas o incorrectas podría contribuir a que tienda a mantenerse más fiel al contexto recuperado, a reconocer con mayor frecuencia cuándo ese contexto es insuficiente para responder, y a emplear la terminología propia de los cursos de ingeniería en la UdelaR.

Técnicas de ajuste fino eficiente como LoRA \cite{hu2022lora} y QLoRA \cite{dettmers2023qlora} permiten adaptar modelos de la escala de Llama 3.1 8B con recursos de cómputo moderados, sin necesidad de reentrenar todos los parámetros del modelo. El proceso requeriría acumular primero datos de calidad a partir de las interacciones reales de los usuarios y luego aplicar el ajuste fino sobre los pares de consulta y respuesta validados. La mejora esperada no es en el recall de la recuperación vectorial, sino en la fidelidad y coherencia de las respuestas generadas a partir de los fragmentos ya recuperados.

El corpus actual indexa únicamente el contenido textual de las transcripciones automáticas. Sin embargo, las clases de OpenFing incluyen diapositivas, diagramas y escritura en pizarra que contienen información complementaria no capturada en el texto transcripto. Extender el índice vectorial para incorporar representaciones de imágenes mediante modelos de visión y lenguaje podría enriquecer el contexto disponible para el asistente, en particular en consultas sobre conceptos que se explican principalmente a través de esquemas visuales o fórmulas escritas en pizarra.

Por último, el pipeline de procesamiento carece actualmente de una suite de tests automatizados que verifique el comportamiento esperado de cada componente de forma aislada y en integración. Añadir tests unitarios para los módulos de chunking, embedding y recuperación, junto con tests de regresión sobre el pipeline RAG completo, reduciría el riesgo de introducir fallas silenciosas al modificar o extender el sistema.


